% ==============================================================================
\chapter{Introduction \& State of the Art}
\label{ch:intro}
% ==============================================================================

\section{Introduction}
Modern engineering design, structural health monitoring, and the realization of Digital Twins require high-fidelity numerical models that can be evaluated in real-time. In structural dynamics, resolving the transient response of complex geometries under varying parameter settings typically relies on the Finite Element Method (FEM). However, high-fidelity FEM simulations of structural components—such as cantilever beams subjected to geometric modifications—demand massive computational resources and time-consuming remeshing, making them unsuitable for interactive design exploration or real-time control loops.

To address these challenges, this thesis investigates the integration of two powerful numerical paradigms: \textbf{Isogeometric Analysis (IGA)} and projection-based \textbf{Reduced Order Modeling (ROM)}. By using Non-Uniform Rational B-Splines (NURBS) as the basis functions for both geometry representation and field approximation, IGA eliminates geometric approximation errors. Concurrently, ROM projects the full-order systems onto low-dimensional subspaces, drastically reducing the size of the systems to be solved.

\section{Conceptual Overview: Everyday Analogies}
Before diving into the mathematical equations of structural dynamics and tensor-product splines, it is helpful to understand the three core concepts of this thesis using simple, everyday analogies:

\begin{itemize}
    \item \textbf{Isogeometric Analysis (Vector Graphics vs. Pixels)}: Traditional simulation (FEM) is like looking at a low-resolution pixelated image. To simulate a curved boundary, standard methods break the smooth curve into thousands of tiny flat segments, losing the original exact shape. Isogeometric Analysis (IGA) is like using vector graphics (like SVG files or Adobe Illustrator designs). The geometry remains mathematically perfect and smooth no matter how much you zoom in, allowing us to change the shape of the beam without generating a new mesh.
    \item \textbf{Reduced Order Modeling (MP3 Audio Compression)}: Simulating a structures' vibration step-by-step usually requires solving millions of equations. Reduced Order Modeling (ROM) is the engineering equivalent of MP3 audio compression. Instead of recording every tiny air pressure wave individually, an MP3 file stores only the principal frequencies (modes) that human ears can hear. Similarly, ROM compresses the millions of structural degrees of freedom down to a handful of primary "shapes of motion" (POD modes), allowing the simulation to solve in milliseconds rather than minutes.
    \item \textbf{Hyper-Reduction (Virtual Sensors)}: When the geometry of the beam changes (e.g., a notch is moved or widened), the computer still has to run calculations over all 3,200 blocks in the mesh, creating a major speed bottleneck. Hyper-reduction (specifically ECSW) is like placing a few dozen virtual "sensors" at critical stress points (such as the notch corners and the loaded tip). By measuring the physical forces only at these key sensor locations and scaling them appropriately, we can accurately predict the behavior of the entire structure while ignoring 99\% of the calculations.
\end{itemize}

Specifically, we consider the parameterized geometry of a cantilever beam with symmetric circular notches, as illustrated in Chapter \ref{ch:theory}. The geometric parameters consist of the notch center location $\mu$ and the notch radius $r$. As these parameters vary, the physical domain $\Omega(\boldsymbol{\mu})$ changes. This variation introduces a vector space incompatibility across different configurations, which is resolved by mapping the equations to a fixed reference configuration $\hat{\Omega}$ (often called "omega hat").

\section{Research Objectives}
The primary objective of this work is to establish a mathematically rigorous, computationally efficient framework for real-time parameterized structural dynamics. The specific goals include:
\begin{enumerate}
    \item \textbf{Exact Geometry Tracking}: Leverage IGA to represent shape parameter variations (such as the notched cantilever beam) exactly using NURBS control point displacements, maintaining a constant topological mesh structure.
    \item \textbf{Dimension Reduction}: Extract an optimal, low-dimensional reduced basis using Proper Orthogonal Decomposition (POD) from snapshots of the transient response.
    \item \textbf{Hyper-Reduction}: Bypass full-mesh assembly for parameter-dependent systems using the Energy-Conserving Sampling and Weighting (ECSW) method to resolve the online computation bottleneck.
    \item \textbf{Real-Time Implementation}: Demonstrate the feasibility of the proposed ROM in a client-side web application, achieving interactive frame rates for real-time visualization of structural responses.
\end{enumerate}

\section{State of the Art}
\subsection{Isogeometric Analysis (IGA)}
Introduced by Hughes et al. in 2005, Isogeometric Analysis (IGA) bridges the gap between Computer-Aided Design (CAD) and Finite Element Analysis (FEA). By utilizing NURBS basis functions for the analysis, IGA offers several key advantages:
\begin{itemize}
    \item Exact geometric representation of curved boundaries without discretization errors.
    \item Higher-order continuity ($C^{p-1}$ for degree $p$), which leads to superior accuracy per degree of freedom compared to classical $C^0$ Lagrange finite elements.
    \item Native support for mesh refinement techniques (knot insertion, degree elevation, and $k$-refinement) while preserving the exact CAD geometry.
\end{itemize}
These properties make IGA particularly well-suited for shape parameterization, as the geometry can be perturbed by simply updating the coordinates of the control points while keeping the knot vectors and element connectivity invariant.

\subsection{Reduced Order Modeling (ROM) for Parameterized Systems}
Projection-based Reduced Order Modeling aims to approximate the high-fidelity Full Order Model (FOM) solution in a low-dimensional subspace. Proper Orthogonal Decomposition (POD), combined with Galerkin projection, is a standard approach.
In a parameterized context, the reduced basis $\boldsymbol{\Phi}$ must capture the system's behavior across the entire parameter domain $\mathcal{D}$. While POD successfully compresses the state space, the online evaluation of the reduced stiffness matrix $\mathbf{K}_{red}(\boldsymbol{\mu})$ remains a bottleneck because the physical domain $\Omega(\boldsymbol{\mu})$ changes, requiring the reconstruction of the full stiffness matrix before projection.

\subsection{Pullback Mappings and Reference Configurations}
To resolve the vector space incompatibility and avoid rebuilding the mesh, the governing equations are pulled back from the parameterized physical domain $\Omega(\boldsymbol{\mu})$ to a fixed, parameter-independent reference configuration $\hat{\Omega}$. By applying a smooth geometric mapping $\boldsymbol{\Psi}(\cdot; \boldsymbol{\mu})$, the integrals are evaluated on $\hat{\Omega}$. Under IGA, the parametric parent space $[0, 1]^2$ serves natively as this reference configuration, providing a mathematically elegant and structurally invariant framework for the assembly process.

\subsection{Hyper-Reduction Methods}
Even when the equations are defined on the reference configuration, evaluating the parameter-dependent integrals at every time step or parameter change requires integration over the entire domain. For nonlinear or parameter-dependent operators, this operation scales with the size of the full-order mesh. To achieve true online efficiency, hyper-reduction methods are employed. The Energy-Conserving Sampling and Weighting (ECSW) method, proposed by Farhat et al., replaces the global domain integration with a sparse, weighted sum over a small subset of elements $\widehat{\Omega} \subset \Omega$. ECSW is particularly robust because it preserves the energy-conserving properties of the underlying physical formulation.

\section{Thesis Outline}
The remainder of this thesis is structured as follows:
\begin{itemize}
    \item \textbf{Chapter \ref{ch:theory} (Theory: Problem Definition)}: Formulates the continuous governing equations of parameterized structural dynamics, defines the geometry configurations, and details the coordinate mapping to the reference configuration.
    \item \textbf{Chapter \ref{ch:discretization} (Discretization: Isogeometric Analysis)}: Details the NURBS mathematical formulation, field approximations, and the assembly of the semi-discrete equations of motion.
    \item \textbf{Chapter \ref{ch:fom} (Simulation: Full Order Model)}: Explains the time integration schemes (Generalized-$\alpha$), geometric nonlinearities, and the generation of snapshot data.
    \item \textbf{Chapter \ref{ch:rom} (Reduced Order Modeling)}: Focuses on POD basis extraction, Galerkin projection, and hyper-reduction via ECSW to achieve real-time simulation speeds.
    \item \textbf{Chapter \ref{ch:digital_twin} (Real-Time Web Application (Digital Twin))}: Introduces the client-side software architecture, reference configuration mapping for real-time visualization, and parameter synchronization.
    \item \textbf{Chapter \ref{ch:conclusion} (Conclusion)}: Summarizes the achievements and outlines future directions for research and optimization.
\end{itemize}
